% Funktionsklassen und Eigenschaften – bank/funktionsklassen-und-eigenschaften/ (zusammenbau v0.4, Heft)
\einheitenkopf[t5]{Funktionsklassen und Eigenschaften}
\setcounter{aufgabe}{50}

% Anlauf (ohne Prüfkennung)
\begin{aufgabe}{Nullstellen – Anlauf}
\begin{teile}
\teil $f(x) = (x + 4)(x - 7)$ – welches Verfahren? \\ \kreuz{Faktoren null setzen} \\ \kreuz{ausklammern} \\ \kreuz{substituieren} \\ \kreuz{Polynomdivision} \\ \kreuz{Lösungsformel}
\teil $f(x) = (x - 2)(x + 5)$ – Nullstellen? x1 = \leerfeld , x2 = \leerfeld
\teil $f(x) = x^3 - 9x$ – Nullstellen?
\end{teile}
\end{aufgabe}

\begin{aufgabe}{Nullstellen – Prüfungsaufgaben}
\begin{teile}
\teil $f(x) = x^3 - 8x^2 + 21x - 18 = (x - 3)(x^2 - 5x + 6)$ – Schnittpunkt mit der y-Achse und Anzahl der Nullstellen \hfill (Abitur 2023 GK)?
\teil $f(x) = x^3 - 3x - 2 = (x + 1)(x^2 - x - 2)$ – Schnittpunkt mit der y-Achse und Anzahl der Nullstellen \hfill (Abitur 2023 GK)?
\end{teile}
\end{aufgabe}

% Anlauf (ohne Prüfkennung)
\begin{aufgabe}{Graph und Term – Anlauf}
\begin{teile}
\teil $f(x) = -2x^4 + x$ – wohin verläuft der Graph für $x \to +\infty$? \\ \kreuz{nach oben} \\ \kreuz{nach unten}
\teil Fülle die Wertetabelle von $f(x) = x^2 - 2$ aus und übertrage die Punkte ins Koordinatensystem.

\wertetabelle{x}{f(x)}{-2,-1,0,1,2}\begin{ksys}[xmin=-3,xmax=3,ymin=-3,ymax=3]\end{ksys}
\teil Lies am Graphen $f(3)$ ab und entscheide mit Begründung, ob $f'(3)$ positiv oder negativ ist.

\begin{ksys}[xmin=-1,xmax=5,ymin=-2,ymax=5,ablesen]\funktion{-0.5*\x^2+2*\x+1}{f}\end{ksys}
\end{teile}
\end{aufgabe}

\begin{aufgabe}{Graph und Term – Prüfungsaufgaben}
\begin{teile}
\teil $f(x) = 3\,\mathrm{cos}(x)$. Drei direkt aufeinanderfolgende Extrempunkte des Graphen, beginnend beim Hochpunkt auf der y-Achse, bilden ein Dreieck. Berechne seinen Flächeninhalt \hfill (Abitur 2026 GK).
\teil $g(x) = 2\,\mathrm{sin}(0{,}5x)$. Die drei ersten Extrempunkte des Graphen rechts der y-Achse bilden ein Dreieck. Berechne seinen Flächeninhalt \hfill (Abitur 2026 GK).
\teil In einer Zeichnung ohne Skalen ist der Graph von $f(x) = x^4 - 2x^3$ dargestellt. Die Nullstelle $2$ liegt $4$ cm rechts vom Ursprung, der Graphenpunkt an der Stelle $-1$ liegt $3$ cm über der x-Achse. Beurteile: Eine Längeneinheit ist auf beiden Achsen gleich lang \hfill (Abitur 2023 GK).
\teil In einer Zeichnung ohne Skalen ist der Graph von $f(x) = 0{,}5x^3 - 2x$ dargestellt. Die Nullstelle $2$ liegt $3$ cm rechts vom Ursprung, der Graphenpunkt an der Stelle $1$ liegt $2{,}25$ cm unter der x-Achse. Beurteile: Eine Längeneinheit ist auf beiden Achsen gleich lang \hfill (Abitur 2023 GK).
\end{teile}
\end{aufgabe}
